% AUTO STUB: replace with a written summary (see writeup/hypothesis-pages/RUBRIC.md). Keep this marker line out of written versions.
\hwhat{\textit{Not yet summarized. The question:} In "pick your own goal" weeks, agents choose and abandon projects freely. Do project abundances follow *neutral cooperative dynamics* (Piñero-style: species, here projects, are equivalent, but joining is frequency-dependent and cooperative)? Signatures: bimodal abundances with a persistent cooperator core, Simpson concentration $\lambda$ near the predicted $\lambda$*($\mu$, N), and boundaries $\mu$\_B (bimodal) and $\mu$\_L (log-series) in the novelty rate $\mu$. Or does plain Hubbell neutral drift (log-series, no cooperation) explain them? Practical payoff: whether a free swarm concentrates its effort on a few shared projects by a predictable law, and what sets how many projects survive.}
\hmodel{From: physics-models/06-neutral-cooperative-dynamics (read its README for the rules, $\lambda$*, $\mu$\_B, $\mu$\_L and the finite-N simulation), with physics-models/10-potts as the categorical-state view. H11 found strong herding (ferromagnetic Potts coupling) on project labels in 11/14 periods, including free weeks \#31 and \#37. That is the cooperative ingredient this model formalizes.}
